\section{SMART goal 2: Investigation of IT systems for DQM metrics}
\subsection*{Goal definition}
\textbf{2) Investigation of IT systems for DQM metrics}

\textbf{WHY:} To define what DQM means in measurable terms.

By the end of the analysis phase, the IT systems and data stores are investigated, and it is defined where quality metrics can be implemented, stored, and evaluated within the DWH LS OPS environment.

Success is measured by a documented system overview that maps the selected metric to its technical source, storage location, and execution context. This is to be validated with IT stakeholders.


\subsection*{Core results}
This goal, while enriched with many supporting results, did not meet the core result of defining the storage and execution context. 
It is important to note that the rule-based DQM system was developed only as a proof of concept with a test case. Therefore, the position of the rule checks was not defined \ref{SolutionDQMEndOf2025}. This is, however, out of scope, as the rule check engine is developed by the DQM team working in parallel on this solution without involvement of the project team.


\subsection*{Supporting results}
\begin{itemize}
\item \textbf{Data lineage description:} The data lineage was described as follows: \textit{1) Physical data source:} Scans are performed on parcels by different scanners; \textit{2) Initial data storage:} The scan events are first saved in AWS S3; \textit{3) Kafka ingestion:} Apache Kafka ingests the events into a topic; \textit{4) Data delivery to DWH:} Apache Kafka is consumed by DWH LS OPS, where the data is saved to the raw layer. This is described in Section \ref{SolutionDataLineage}.

\item \textbf{Reporting pipeline description:} The reporting pipeline of DWH LS OPS is defined as follows: the \textbf{storage layer} using Amazon S3, followed by the \textbf{query and processing layer} using Amazon Athena and SAP Datasphere (DSP), and finally the \textbf{consumption layer} using SAP Analytics Cloud (SAC). These steps are detailed in Section \ref{SolutionReportingPipeline}.

\item \textbf{Placement of data quality checks:} Data quality checks are most effectively placed on the processing and transformation layer. This allows both meaningful and effective validation. Earlier checks on the data lineage and reporting pipeline are limited to global completeness, while later checks on the reporting layer are too late to enable corrective action. This is discussed in Section \ref{ResearchDataQChecksOnDataLifecycle}.

\item \textbf{Results of DQM checks:} The DQM team defined that the results of rule checks are to be stored in a rule-check database. However, this solution was not yet installed in DWH LS OPS at the time of this analysis. This is described in Section \ref{SolutionDQMTeamHighLevelTarget}.
\end{itemize}


\subsection*{What was not achieved?}

\begin{itemize}
    
\item \textbf{IT stakeholder confirmation:} Formal confirmation of the data lineage and the reporting pipeline was not obtained from IT stakeholders.

\item \textbf{DQM check storage location:} The location where DQM check results are to be stored was not defined.

\end{itemize}


\subsection*{Goal verdict}
This SMART goal was \textbf{partially achieved}, as two targets were not met.

\textbf{Why:} First, the DQM check storage location was not defined; however, this is out of scope because responsibility for the application lies with the DQM team. Second, formal confirmation with IT stakeholders regarding the data lineage and reporting pipeline was not conducted.


\subsection*{Learning and outlook for future work}
Regarding the missing alignment with IT stakeholders, the data lineage was initially considered general knowledge within the team and therefore not documented first. Once it was available, discussion with IT stakeholders was no longer possible due to time constraints. Second, the position for saving DQM results was out of scope.

For the first point, aspects that appear obvious should be documented earlier when confirmation is required. Regarding out-of-scope topics, communication with external teams needs to be more proactive in order to identify earlier when deliverables are delayed or when priorities change.